\chapter{Visão de Cenários (+1 View -- Casos de Uso Críticos)}
\label{cap_visao_cenarios}

No Modelo 4+1 de Kruchten, a Visão de Cenários -- representada pelo ``+1'' da fórmula -- desempenha o papel essencial de unificador e validador da arquitetura de software \cite{kruchten1995}. Essa visão seleciona os casos de uso mais críticos e arquiteturalmente significativos do sistema, demonstrando como os elementos projetados nas demais visões (classes na Visão Lógica, processos na Visão de Processos, pacotes na Visão de Desenvolvimento e nós na Visão Física) colaboram harmoniosamente para materializar o comportamento exigido \cite{bass2021,pressman2021}.

% ===========================================================================
\section{Diagrama Geral de Casos de Uso Arquiteturais}
\label{sec_diagrama_casos_uso_arq}
% ===========================================================================

A \autoref{fig_usecase_arquitetura} apresenta o Diagrama Geral de Casos de Uso do SIGEA-GTT na notação UML 2.5, sintetizando as interações dos três perfis do sistema (\texttt{ADMIN}, \texttt{PROFESSOR} e \texttt{STUDENT}) que exercitam os pontos nevrálgicos da arquitetura.

\begin{figure}[p]
\centering
\caption{Diagrama de Casos de Uso Arquiteturais do SIGEA-GTT (UML 2.5).}
\label{fig_usecase_arquitetura}
\resizebox{0.92\textwidth}{!}{%
\begin{tikzpicture}[
  >=Stealth,
  uc/.style={
    ellipse, draw=blue!70!black, fill=blue!5, thick,
    font=\tiny, align=center, minimum width=3.3cm, minimum height=0.75cm
  },
  bdr/.style={
    rectangle, draw=blue!40, dashed, thick, rounded corners=6pt,
    fill=blue!2
  },
  act/.style={font=\small\bfseries, draw=blue!80!black, fill=white, rounded corners=2pt, inner sep=4pt}
]

% Sistema Boundary
\node[bdr, minimum width=9.5cm, minimum height=15.5cm]
     (sys) at (5, -6) {};
\node[above, font=\small\bfseries, blue!80!black] at (sys.north) {SIGEA-GTT -- Fronteira Arquitetural do Sistema};

% Atores
\node[act] (admin) at (-2,  -1.5) {ADMIN};
\node[act] (prof)  at (-2,  -7.5) {PROFESSOR};
\node[act] (aluno) at (12,  -8.5) {STUDENT};

% Casos de uso comuns
\node[uc] (login)   at (5,   0.5) {UC01: Autenticar via JWT\\(RF01, RF02, RNF03)};
\node[uc] (senha)   at (5,  -0.8) {UC02: Trocar Senha Obrigatória\\(RF02, RF03)};
\node[uc] (catalogo)at (5,  -2.1) {UC03: Consultar Catálogo GTT\\e Guia MERP (RF08, RF09)};

% Casos de uso Admin
\node[uc] (uc_users)   at (3, -3.8) {UC04: Gerenciar Usuários e Matrículas\\(RF04, RF05, RF12)};
\node[uc] (uc_modgtt)  at (3, -5.2) {UC05: CRUD Módulos e Gatilhos\\(RF06, RF07)};
\node[uc] (uc_turmas)  at (3, -6.6) {UC06: Administrar Turmas\\(RF10, RF11, RF13)};
\node[uc] (uc_dash_a)  at (3, -8.0) {UC07: Monitorar Métricas Gerais\\(RF25, RNF02)};

% Casos de uso Professor
\node[uc] (uc_atv)     at (7, -3.8) {UC08: Elaborar Atividade com\\Prontuário JSONB (RF14, RF15)};
\node[uc] (uc_corr)    at (7, -5.2) {UC09: Avaliar Submissão e\\Emitir Parecer (RF24)};
\node[uc] (uc_dash_p)  at (7, -6.6) {UC10: Analisar Indicadores da\\Turma (RF25, RNF02)};

% Casos de uso Aluno
\node[uc] (uc_cron)    at (7, -8.2) {UC11: Auditar Prontuário sob\\Cronômetro 20 min (RF16)};
\node[uc] (uc_gatid)   at (7, -9.6) {UC12: Identificar Gatilhos\\e Danos MERP (RF17)};
\node[uc] (uc_ferrq)   at (7,-11.0) {UC13: Aplicar 6 Ferramentas\\da Qualidade (RF18 a RF22)};
\node[uc] (uc_sub)     at (7,-12.4) {UC14: Submeter Resolução\\Definitiva (RF23)};

% Relacionamentos Admin
\draw[->] (admin) -- (login);
\draw[->] (admin) -- (uc_users);
\draw[->] (admin) -- (uc_modgtt);
\draw[->] (admin) -- (uc_turmas);
\draw[->] (admin) -- (uc_dash_a);

% Relacionamentos Professor
\draw[->] (prof) -- (login);
\draw[->] (prof) -- (catalogo);
\draw[->] (prof) -- (uc_atv);
\draw[->] (prof) -- (uc_corr);
\draw[->] (prof) -- (uc_dash_p);

% Relacionamentos Aluno
\draw[->] (aluno) -- (login);
\draw[->] (aluno) -- (catalogo);
\draw[->] (aluno) -- (uc_cron);
\draw[->] (aluno) -- (uc_gatid);
\draw[->] (aluno) -- (uc_ferrq);
\draw[->] (aluno) -- (uc_sub);

% Include
\draw[->, dashed] (login) -- node[right, font=\tiny]{include} (senha);
\draw[->, dashed] (uc_cron) -- (uc_gatid);
\draw[->, dashed] (uc_gatid) -- (uc_ferrq);
\draw[->, dashed] (uc_ferrq) -- (uc_sub);

\end{tikzpicture}%
}
\fonte{Elaborado pelo autor com TikZ (2026).}
\end{figure}

% ===========================================================================
\section{Cenários Arquiteturais Críticos}
\label{sec_cenarios_criticos}
% ===========================================================================

Quatro cenários foram identificados como condutores prioritários da tomada de decisões arquiteturais do SIGEA-GTT.

% ---------------------------------------------------------------------------
\subsection{Cenário 1: Autenticação Institucional e Troca de Senha Provisória}
\label{subsec_cenario_auth}
% ---------------------------------------------------------------------------
\begin{itemize}
    \item \textbf{Estímulo}: O usuário informa e-mail \texttt{@academico.ufs.br} e senha provisória emitida pelo administrador;
    \item \textbf{Requisitos Vinculados}: \texttt{RF01}, \texttt{RF02}, \texttt{RF03}, \texttt{RNF03};
    \item \textbf{Implicações Arquiteturais}:
    \begin{enumerate}
        \item O backend valida as credenciais contra a tabela \texttt{users} utilizando a função criptográfica de derivação BCrypt com fator de custo 12;
        \item O serviço de autenticação emite um JWT com assinatura HMAC-SHA256 e insere no \textit{payload} a reivindicação (\textit{claim}) booleana \texttt{primeiroAcesso: true};
        \item O interceptor HTTP do Angular insere o cabeçalho \texttt{Authorization: Bearer <token>} em todas as requisições subsequentes;
        \item As guardas de rota do Angular (\texttt{FirstAccessGuard}) interceptam a navegação e forçam o redirecionamento compulsoriamente para a tela de redefinição de senha, bloqueando qualquer acesso a recursos acadêmicos enquanto a nova credencial não for confirmada e persistida.
    \end{enumerate}
\end{itemize}

% ---------------------------------------------------------------------------
\subsection{Cenário 2: Auditoria sob Cronômetro de 20 Minutos e Prontuário JSONB}
\label{subsec_cenario_auditoria}
% ---------------------------------------------------------------------------
\begin{itemize}
    \item \textbf{Estímulo}: O estudante inicia a resolução de uma atividade prática vinculada à sua turma;
    \item \textbf{Requisitos Vinculados}: \texttt{RF15}, \texttt{RF16}, \texttt{RF17}, \texttt{RNF01}, \texttt{RNF06}, \texttt{RNF07};
    \item \textbf{Implicações Arquiteturais}:
    \begin{enumerate}
        \item O backend extrai da entidade \texttt{Atividade} o prontuário clínico simulado persistido no PostgreSQL como tipo nativo \texttt{JSONB} (contendo dados de admissão, evoluções diárias, prescrições e exames laboratoriais);
        \item O frontend Angular 18 faz o \textit{parse} reativo e organiza o prontuário em abas temáticas de alta densidade visual calibradas com a paleta cromática clínica hospitalar suave (\texttt{RNF06});
        \item Um serviço reativo no frontend inicia o cronômetro regressivo de 20 minutos (1200 segundos), utilizando um \textit{Signal} com temporização a cada 1000 ms, disparando avisos cromáticos quando restam 5 minutos (âmbar) e 1 minuto (carmim);
        \item O discente seleciona os gatilhos clínicos identificados a partir do catálogo e especifica se houve dano real (classificação NCC MERP de E a I).
    \end{enumerate}
\end{itemize}

% ---------------------------------------------------------------------------
\subsection{Cenário 3: Aplicação Integrada das Seis Ferramentas da Qualidade}
\label{subsec_cenario_qualidade}
% ---------------------------------------------------------------------------
\begin{itemize}
    \item \textbf{Estímulo}: O estudante investiga as causas-raízes e propõe contramedidas hospitalares preventivas;
    \item \textbf{Requisitos Vinculados}: \texttt{RF18}, \texttt{RF19}, \texttt{RF20}, \texttt{RF21}, \texttt{RF22};
    \item \textbf{Implicações Arquiteturais}:
    \begin{enumerate}
        \item O componente interativo do Diagrama de Ishikawa 6M permite a inclusão e exclusão de nós causais dinâmicos em torno do problema central;
        \item A Matriz GUT calcula reativamente o produto $Score = G \times U \times T$ via \textit{computed signals} do Angular a cada alteração de nota (1 a 5), ordenando instantaneamente a lista de prioridades;
        \item As matrizes 5W2H, PDCA, SWOT e Brainstorming são agregadas em um documento estruturado persistido no campo \texttt{submissoes.ferramentas\_qualidade} em formato \texttt{JSONB}.
    \end{enumerate}
\end{itemize}

% ---------------------------------------------------------------------------
\subsection{Cenário 4: Avaliação Docente e Consolidação de Métricas Epidemiológicas}
\label{subsec_cenario_avaliacao}
% ---------------------------------------------------------------------------
\begin{itemize}
    \item \textbf{Estímulo}: O professor titular examina as submissões entregues e atribui a nota com parecer formativo;
    \item \textbf{Requisitos Vinculados}: \texttt{RF24}, \texttt{RF25}, \texttt{RNF02};
    \item \textbf{Implicações Arquiteturais}:
    \begin{enumerate}
        \item O backend valida que o parecer pedagógico formativo seja obrigatório e que a nota se situe no intervalo $[0,00; 10,00]$;
        \item Ao atualizar o status da submissão para \texttt{CORRIGIDO}, o serviço de indicadores recomputa em tempo de execução os três indicadores epidemiológicos clássicos do IHI:
        \begin{align}
            TI_{1000} &= \left( \frac{\sum \text{Eventos Adversos}}{\sum \text{Dias de Internação}} \right) \times 1.000 \\
            TI_{100}  &= \left( \frac{\sum \text{Eventos Adversos}}{\text{Casos Auditados}} \right) \times 100 \\
            P_{EA}    &= \left( \frac{\text{Casos com } \ge 1 \text{ EA}}{\text{Casos Auditados}} \right) \times 100
        \end{align}
        \item A consulta agregada no PostgreSQL utiliza índices e projeções DTO do MapStruct, garantindo que o tempo de resposta permaneça rigorosamente abaixo do limiar de 200 ms (\texttt{RNF02}).
    \end{enumerate}
\end{itemize}
